\documentclass[a4paper,12pt]{article}

\usepackage[T2A]{fontenc}
\usepackage[utf8]{inputenc}
\usepackage[russian]{babel}
\usepackage{amsmath,amssymb,amsthm,mathtools}
\usepackage{helvet}
\renewcommand{\familydefault}{\sfdefault}
\usepackage{icomma}
\usepackage{geometry}
\geometry{left=20mm,right=20mm,top=20mm,bottom=22mm}
\usepackage{microtype}
\usepackage{tikz}
\usetikzlibrary{calc}
\usetikzlibrary{arrows.meta,positioning,shapes.geometric}
\usepackage{tkz-euclide}
\usepackage{pgfplots}
\pgfplotsset{compat=1.18}
\usepackage{tabularx,booktabs,longtable}
\usepackage{enumitem}
\usepackage{hyperref}
\usepackage{parskip}
\usepackage{fancyhdr}
\usepackage{setspace}
\usepackage{xcolor}

\definecolor{accent}{HTML}{008080}
\definecolor{darkaccent}{HTML}{004D4D}
\definecolor{textgray}{HTML}{333333}
\definecolor{softgray}{HTML}{F2F4F4}

\hypersetup{
    colorlinks=true,
    linkcolor=darkaccent,
    urlcolor=darkaccent
}

\onehalfspacing
\setlength{\parindent}{0pt}
\setlength{\headheight}{25pt}
\setlist[itemize]{leftmargin=6mm,itemsep=1mm,topsep=1mm}
\setlist[enumerate]{leftmargin=7mm,itemsep=2mm,topsep=1mm}
\renewcommand{\arraystretch}{1.35}

\pagestyle{fancy}
\fancyhf{}
\fancyhead[L]{\small\textcolor{textgray}{ОГЭ по математике}}
\fancyhead[R]{\small\textcolor{textgray}{София Кальней}}
\fancyfoot[C]{\small\thepage}
\renewcommand{\headrulewidth}{0.4pt}
\renewcommand{\headrule}{
    \hbox to\headwidth{
        \color{accent}
        \leaders\hrule height \headrulewidth\hfill
    }
}

\tkzSetUpLine[line width=1pt,color=accent]
\tkzSetUpPoint[color=accent,fill=accent,size=3]
\tkzSetUpLabel[color=black,font=\small]

\tikzset{
    flowstart/.style={
        ellipse,
        draw=darkaccent,
        line width=0.8pt,
        align=center,
        minimum width=35mm,
        minimum height=9mm,
        font=\small
    },
    flowaction/.style={
        rounded corners=2mm,
        rectangle,
        draw=accent,
        line width=0.8pt,
        align=center,
        text width=105mm,
        minimum height=10mm,
        font=\small,
        inner sep=4pt
    },
    flowdecision/.style={
        diamond,
        aspect=2.4,
        draw=darkaccent,
        line width=0.8pt,
        align=center,
        text width=58mm,
        inner sep=2pt,
        font=\small
    },
    flowarrow/.style={
        -{Stealth[length=3mm]},
        line width=0.8pt,
        draw=textgray
    }
}

\newcommand{\topicline}{
    \textcolor{accent}{\rule{\linewidth}{0.6pt}}
}

\newcommand{\checkitem}[1]{
    \item[$\square$] #1
}

\begin{document}

% =========================================================
% ТИТУЛЬНЫЙ ЛИСТ
% =========================================================

\begin{titlepage}
\thispagestyle{empty}
\centering

\vspace*{22mm}

{\Large\textcolor{darkaccent}{\textbf{ОГЭ по математике}}\par}

\vspace{12mm}

{\Huge\textbf{Чек-лист}\par}

\vspace{5mm}

{\LARGE\textcolor{accent}{
Повторение задач повышенной сложности
}\par}

\vspace{4mm}

{\large
Алгебра, графики и геометрические конструкции
\par}

\vspace{18mm}

{\large\textbf{Ученица:} София Кальней\par}

\vspace{4mm}

{\large\textbf{Дата:} \today\par}

\vfill

{\large Лёвин Артём Александрович\par}
{\normalsize эксклюзивно для Софии Кальней\par}

\vspace{13mm}

\begin{tikzpicture}[x=1cm,y=1cm]
    \draw[
        accent,
        line width=0.9pt
    ]
    (0,0)
    .. controls (3,0.8) and (6,-0.8) ..
    (9,0)
    .. controls (11,0.45) and (13,-0.35) ..
    (15,0.15);
\end{tikzpicture}

\vspace{5mm}

\end{titlepage}

% =========================================================
% 1. КАРТА ЗАНЯТИЯ
% =========================================================

\section*{1. Карта занятия}
\topicline

Цель повторения --- быстро узнавать тип задачи, выбирать устойчивый
метод и сохранять ограничения исходного условия на каждом этапе
преобразований.

\begin{itemize}
    \checkitem{
        Рациональные неравенства:
        нули знаменателя, знаки на интервалах,
        строгие и нестрогие границы.
    }

    \checkitem{
        Текстовые задачи на работу:
        производительность, время, объём работы,
        перевод условия в уравнение.
    }

    \checkitem{
        Рациональные функции:
        разложение на множители,
        сокращение и сохранение выколотых точек.
    }

    \checkitem{
        Параметрические прямые $y=kx$:
        параллельность и прохождение
        через выколотые точки.
    }

    \checkitem{
        Системы уравнений:
        подстановка, замена $t=y^2$,
        альтернативное использование $(x+y)^2$.
    }

    \checkitem{
        Графики с модулем:
        раскрытие по случаям
        и отбор нужных частей парабол.
    }

    \checkitem{
        Прямоугольный треугольник:
        теорема Пифагора и высота к гипотенузе
        через равенство площадей.
    }

    \checkitem{
        Трапеция:
        диагональ, средняя линия,
        подобие треугольников
        и дополнительные построения.
    }
\end{itemize}

% =========================================================
% 2. РАЦИОНАЛЬНОЕ НЕРАВЕНСТВО
% =========================================================

\section*{2. Рациональное неравенство: знак дроби}
\topicline

Рассмотрим типовую конструкцию

\[
\frac{-14}{x^2+5x-14}\le 0.
\]

Знаменатель раскладывается по теореме Виета:

\[
x^2+5x-14=(x-2)(x+7).
\]

Поэтому

\[
x\ne2,
\qquad
x\ne-7.
\]

Числитель $-14$ всегда отрицателен и никогда не обращается в нуль.
Следовательно, дробь неположительна именно там, где знаменатель
положителен:

\[
(x-2)(x+7)>0.
\]

Отсюда

\[
\boxed{
x\in(-\infty,-7)\cup(2,+\infty)
}.
\]

\subsection*{Числовая прямая}

\begin{center}
\begin{tikzpicture}[x=1cm,y=1cm]

    \draw[
        -{Stealth[length=2.5mm]},
        line width=0.8pt
    ]
    (-6.5,0)--(6.5,0);

    \draw
    (-2.6,0.12)--(-2.6,-0.12)
    node[below=3pt] {$-7$};

    \draw
    (2.6,0.12)--(2.6,-0.12)
    node[below=3pt] {$2$};

    \draw[
        fill=white,
        line width=0.9pt
    ]
    (-2.6,0) circle (2.2pt);

    \draw[
        fill=white,
        line width=0.9pt
    ]
    (2.6,0) circle (2.2pt);

    \node at (-4.5,0.55) {$-$};
    \node at (0,0.55) {$+$};
    \node at (4.5,0.55) {$-$};

    \node[
        align=center,
        text width=35mm
    ]
    at (0,-1.0)
    {знак всей дроби};

\end{tikzpicture}
\end{center}

\textbf{Контрольная мысль.}
Точки $x=-7$ и $x=2$ остаются выколотыми при любом знаке
неравенства, поскольку знаменатель исходной дроби в них равен нулю.

% =========================================================
% БЛОК-СХЕМА 1
% =========================================================

\clearpage

\section*{Алгоритм: рациональное неравенство методом интервалов}

\vspace{8mm}

\begin{center}
\begin{tikzpicture}[node distance=8mm]

    \node[flowstart] (s)
    {Начало};

    \node[
        flowaction,
        below=of s
    ] (a1)
    {
        Перенесите всё в одну сторону
        и представьте выражение одной дробью
    };

    \node[
        flowaction,
        below=of a1
    ] (a2)
    {
        Разложите числитель
        и знаменатель на множители
    };

    \node[
        flowaction,
        below=of a2
    ] (a3)
    {
        Найдите критические точки:
        нули числителя и нули знаменателя
    };

    \node[
        flowaction,
        below=of a3
    ] (a4)
    {
        Отметьте точки на числовой прямой;
        нули знаменателя всегда выколоты
    };

    \node[
        flowaction,
        below=of a4
    ] (a5)
    {
        Определите знак выражения
        на каждом интервале
    };

    \node[
        flowdecision,
        below=10mm of a5
    ] (d1)
    {
        Нужна граница при знаке
        $\le$ или $\ge$?
    };

    \node[
        flowaction,
        below left=12mm and 10mm of d1,
        text width=42mm
    ] (yes)
    {
        Включите только нули числителя,
        если выражение определено
    };

    \node[
        flowaction,
        below right=12mm and 10mm of d1,
        text width=42mm
    ] (no)
    {
        При строгом знаке
        все критические точки
        остаются вне ответа
    };

    \node[
        flowstart,
        below=20mm of d1
    ] (f)
    {
        Запишите объединение
        подходящих интервалов
    };

    \draw[flowarrow] (s)--(a1);
    \draw[flowarrow] (a1)--(a2);
    \draw[flowarrow] (a2)--(a3);
    \draw[flowarrow] (a3)--(a4);
    \draw[flowarrow] (a4)--(a5);
    \draw[flowarrow] (a5)--(d1);

    \draw[flowarrow]
    (d1)--node[above left,font=\small]{да}(yes);

    \draw[flowarrow]
    (d1)--node[above right,font=\small]{нет}(no);

    \draw[flowarrow]
    (yes.south) |- (f.west);

    \draw[flowarrow]
    (no.south) |- (f.east);

\end{tikzpicture}
\end{center}

\clearpage

% =========================================================
% 3. ТЕКСТОВЫЕ ЗАДАЧИ
% =========================================================

\section*{3. Текстовые задачи на работу}
\topicline

Основная модель:

\[
A=pt,
\]

где $A$ --- выполненная работа,
$p$ --- производительность,
$t$ --- время.

\begin{center}
\begin{tabularx}{\linewidth}{@{}lXXX@{}}
\toprule
Исполнитель
&
Производительность
&
Время
&
Работа
\\
\midrule
Первый
&
$p_1$
&
$t_1$
&
$p_1t_1$
\\
Второй
&
$p_2$
&
$t_2$
&
$p_2t_2$
\\
\bottomrule
\end{tabularx}
\end{center}

\textbf{Практический порядок действий:}

\begin{enumerate}
    \item
    Введите переменную для величины,
    которую проще всего выразить из условия.

    \item
    Сразу подпишите единицы измерения.

    \item
    Заполните одну строку таблицы целиком,
    затем вторую.

    \item
    Составьте уравнение по условию
    о времени, производительности
    или общем объёме.

    \item
    До громоздких преобразований
    сократите общие числовые множители,
    если это возможно.
\end{enumerate}

\textbf{Типичная ошибка:}
складывать времена или производительности без смыслового основания.
Сначала формулируется связь из текста задачи,
затем появляется знак равенства.

% =========================================================
% 4. РАЦИОНАЛЬНАЯ ФУНКЦИЯ
% =========================================================

\section*{4. Рациональная функция: сокращение и выколотые точки}
\topicline

Рассмотрим функцию

\[
y=
\frac{(x+1)(x^2-4)}
     {x^2-x-2}.
\]

Факторизуем:

\[
x^2-4=(x-2)(x+2),
\]

\[
x^2-x-2=(x-2)(x+1).
\]

Исходная область допустимых значений:

\[
\boxed{
x\ne2,
\qquad
x\ne-1
}.
\]

После сокращения получаем

\[
y=x+2,
\]

но ограничения исходной функции сохраняются.

Поэтому график --- прямая

\[
y=x+2
\]

с двумя выколотыми точками:

\[
(2,4),
\qquad
(-1,1).
\]

\begin{center}
\begin{tikzpicture}
\begin{axis}[
    width=0.82\linewidth,
    height=0.48\textheight,
    axis lines=middle,
    xmin=-5,
    xmax=5,
    ymin=-3,
    ymax=7,
    xtick={-4,-3,-2,-1,1,2,3,4},
    ytick={-2,-1,1,2,3,4,5,6},
    unit vector ratio=1 1,
    clip=false,
    xlabel={$x$},
    ylabel={$y$},
    samples=220
]

\addplot[
    accent,
    line width=1.2pt,
    domain=-5:5
]
{x+2};

\addplot[
    only marks,
    mark=o,
    mark size=3pt,
    line width=1pt,
    accent,
    fill=white
]
coordinates {
    (-1,1)
    (2,4)
};

\node[
    anchor=west
]
at (axis cs:2.15,4.25)
{$(2,4)$};

\node[
    anchor=east
]
at (axis cs:-1.15,0.75)
{$(-1,1)$};

\end{axis}
\end{tikzpicture}
\end{center}

\subsection*{Когда прямая $y=kx$ не имеет общих точек с графиком?}

Обычная точка пересечения с прямой $y=x+2$ удовлетворяет

\[
kx=x+2.
\]

Общих точек не будет в трёх случаях.

\begin{enumerate}
    \item
    $k=1$:
    прямые $y=x$ и $y=x+2$ параллельны.

    \item
    Единственная точка пересечения
    формально совпадает с выколотой точкой $(2,4)$:

    \[
    k=\frac42=2.
    \]

    \item
    Единственная точка пересечения
    формально совпадает с выколотой точкой $(-1,1)$:

    \[
    k=\frac{1}{-1}=-1.
    \]
\end{enumerate}

Итак,

\[
\boxed{
k\in\{-1,1,2\}
}.
\]

% =========================================================
% БЛОК-СХЕМА 2
% =========================================================

\clearpage

\section*{Алгоритм: построение рационального графика после сокращения}

\vspace{8mm}

\begin{center}
\begin{tikzpicture}[node distance=8mm]

    \node[flowstart] (s)
    {Исходная рациональная функция};

    \node[
        flowaction,
        below=of s
    ] (a1)
    {
        Сначала найдите ОДЗ
        исходного выражения
    };

    \node[
        flowaction,
        below=of a1
    ] (a2)
    {
        Разложите числитель
        и знаменатель на множители
    };

    \node[
        flowdecision,
        below=10mm of a2
    ] (d1)
    {
        Есть общие множители?
    };

    \node[
        flowaction,
        below left=13mm and 10mm of d1,
        text width=42mm
    ] (yes)
    {
        Сократите их,
        но сохраните запрещённые значения $x$
    };

    \node[
        flowaction,
        below right=13mm and 10mm of d1,
        text width=42mm
    ] (no)
    {
        Стройте исходную функцию
        с учётом её ОДЗ
    };

    \node[
        flowaction,
        below=34mm of d1
    ] (a3)
    {
        Постройте упрощённый график
    };

    \node[
        flowaction,
        below=of a3
    ] (a4)
    {
        В запрещённых значениях вычислите координаты
        соответствующих точек упрощённого графика
        и сделайте их выколотыми
    };

    \node[
        flowstart,
        below=of a4
    ] (f)
    {
        График полностью соответствует
        исходной функции
    };

    \draw[flowarrow] (s)--(a1);
    \draw[flowarrow] (a1)--(a2);
    \draw[flowarrow] (a2)--(d1);

    \draw[flowarrow]
    (d1)--node[above left,font=\small]{да}(yes);

    \draw[flowarrow]
    (d1)--node[above right,font=\small]{нет}(no);

    \draw[flowarrow]
    (yes.south) |- (a3.west);

    \draw[flowarrow]
    (no.south) |- (a3.east);

    \draw[flowarrow]
    (a3)--(a4);

    \draw[flowarrow]
    (a4)--(f);

\end{tikzpicture}
\end{center}

\clearpage

% =========================================================
% 5. СИСТЕМА УРАВНЕНИЙ
% =========================================================

\section*{5. Система уравнений: два способа}
\topicline

Рассмотрим систему

\[
\begin{cases}
x^2+y^2=40,\\
xy=-12.
\end{cases}
\]

\subsection*{Способ 1. Подстановка и замена}

Из второго уравнения при $y\ne0$:

\[
x=-\frac{12}{y}.
\]

Подставляем в первое:

\[
\frac{144}{y^2}+y^2=40.
\]

Поскольку повторяется выражение $y^2$, вводим замену

\[
t=y^2,
\qquad
t>0.
\]

Получаем

\[
\frac{144}{t}+t=40.
\]

Умножаем на $t$:

\[
144+t^2=40t.
\]

Следовательно,

\[
t^2-40t+144=0.
\]

Дискриминант:

\[
D=
40^2-4\cdot144
=
1600-576
=
1024
=
32^2.
\]

Отсюда

\[
t_1=
\frac{40+32}{2}
=
36,
\]

\[
t_2=
\frac{40-32}{2}
=
4.
\]

Значит,

\[
y^2=36
\quad\Longrightarrow\quad
y=\pm6,
\]

или

\[
y^2=4
\quad\Longrightarrow\quad
y=\pm2.
\]

По условию

\[
xy=-12.
\]

Получаем четыре пары:

\[
\boxed{
(-2,6),\;
(2,-6),\;
(-6,2),\;
(6,-2)
}.
\]

\subsection*{Способ 2. Формула квадрата суммы}

Умножим второе уравнение на $2$:

\[
2xy=-24.
\]

Складываем его с равенством

\[
x^2+y^2=40.
\]

Получаем

\[
x^2+2xy+y^2=16.
\]

Левая часть --- полный квадрат:

\[
(x+y)^2=16.
\]

Следовательно,

\[
x+y=4
\]

или

\[
x+y=-4.
\]

В первом случае

\[
x=4-y,
\]

во втором

\[
x=-4-y.
\]

После подстановки в $xy=-12$
каждый случай сводится к обычному квадратному уравнению.

\textbf{Полезная идея.}
Если в системе известны $x^2+y^2$ и $xy$,
стоит сразу проверить формулы

\[
(x+y)^2
=
x^2+2xy+y^2,
\]

\[
(x-y)^2
=
x^2-2xy+y^2.
\]

% =========================================================
% 6. ГРАФИК С МОДУЛЕМ
% =========================================================

\section*{6. График с модулем: две ветви парабол}
\topicline

Рассмотрим функцию

\[
y=x^2+4|x|-5.
\]

Раскрываем модуль по знаку $x$:

\[
y=
\begin{cases}
x^2+4x-5, & x\ge0,\\
x^2-4x-5, & x<0.
\end{cases}
\]

Каждая формула задаёт параболу,
однако в итоговом графике остаётся только та её часть,
которая соответствует своему условию на $x$.

\subsection*{Первая парабола}

\[
y=x^2+4x-5.
\]

Координата вершины:

\[
x_0=
-\frac{4}{2}
=
-2.
\]

Значение функции:

\[
y(-2)
=
4-8-5
=
-9.
\]

Полная парабола имеет вершину

\[
(-2,-9).
\]

Нули:

\[
x^2+4x-5=0.
\]

По теореме Виета

\[
x_1=-5,
\qquad
x_2=1.
\]

Но для этой ветви действует условие

\[
x\ge0.
\]

Следовательно, из двух корней сохраняется только

\[
x=1.
\]

\subsection*{Вторая парабола}

\[
y=x^2-4x-5.
\]

Корни:

\[
x_1=-1,
\qquad
x_2=5.
\]

Для этой ветви действует условие

\[
x<0.
\]

Поэтому сохраняется только корень

\[
x=-1.
\]

\begin{center}
\begin{tikzpicture}
\begin{axis}[
    width=0.82\linewidth,
    height=0.47\textheight,
    axis lines=middle,
    xmin=-4.3,
    xmax=4.3,
    ymin=-6,
    ymax=15,
    xtick={-4,-3,-2,-1,1,2,3,4},
    ytick={-5,5,10,15},
    clip=false,
    xlabel={$x$},
    ylabel={$y$},
    samples=220
]

\addplot[
    accent,
    line width=1.2pt,
    domain=0:4.3
]
{x^2+4*x-5};

\addplot[
    accent,
    line width=1.2pt,
    domain=-4.3:0
]
{x^2-4*x-5};

\addplot[
    only marks,
    mark=*,
    mark size=2.4pt,
    accent
]
coordinates {
    (0,-5)
    (-1,0)
    (1,0)
};

\end{axis}
\end{tikzpicture}
\end{center}

Из рисунка видно:
горизонтальная прямая может иметь
не более двух общих точек с графиком.

Следовательно,

\[
\boxed{2}.
\]

\textbf{Полезная привычка.}
При построении полной параболы выбирайте точки,
симметричные относительно её оси.
После этого оставляйте только ту часть,
которая удовлетворяет условию соответствующей ветви.

% =========================================================
% 7. ПРЯМОУГОЛЬНЫЙ ТРЕУГОЛЬНИК
% =========================================================

\section*{7. Прямоугольный треугольник: высота к гипотенузе}
\topicline

Пусть в прямоугольном треугольнике
гипотенуза равна $74$,
а один катет равен $24$.

По теореме Пифагора второй катет:

\[
b=
\sqrt{74^2-24^2}.
\]

Используем разность квадратов:

\[
74^2-24^2
=
(74-24)(74+24).
\]

Следовательно,

\[
b=
\sqrt{50\cdot98}.
\]

Замечаем удобную перегруппировку:

\[
50\cdot98
=
100\cdot49.
\]

Поэтому

\[
b=
\sqrt{100\cdot49}
=
70.
\]

Теперь найдём высоту $h$,
проведённую к гипотенузе.

Площадь прямоугольного треугольника
через катеты:

\[
S=
\frac12\cdot24\cdot70.
\]

Та же площадь через гипотенузу и высоту:

\[
S=
\frac12\cdot74\cdot h.
\]

Приравниваем:

\[
\frac12\cdot24\cdot70
=
\frac12\cdot74\cdot h.
\]

Отсюда

\[
h=
\frac{24\cdot70}{74}
=
\boxed{\frac{840}{37}}.
\]

\begin{center}
\begin{tikzpicture}[scale=0.9]

    \tkzDefPoint(0,0){A}
    \tkzDefPoint(7,0){B}
    \tkzDefPoint(0,2.4){C}
    \tkzDefPoint(0.736,2.148){H}

    \tkzDrawPolygon(A,B,C)
    \tkzDrawSegment(A,H)

    \tkzDrawPoints(A,B,C,H)

    \tkzLabelPoints[below](A,B)
    \tkzLabelPoints[above](C,H)

    \tkzMarkRightAngle[size=0.28](B,A,C)
    \tkzMarkRightAngle[size=0.28](A,H,C)

    \tkzLabelSegment[left](A,C){$24$}
    \tkzLabelSegment[above right](B,C){$74$}
    \tkzLabelSegment[above left](A,H){$h$}

\end{tikzpicture}
\end{center}

\textbf{Идея, которую стоит помнить.}
Высота очень часто появляется через площадь.
Если одна и та же площадь вычисляется двумя способами,
неизвестная высота извлекается из равенства формул.

% =========================================================
% 8. ДОКАЗАТЕЛЬСТВО О ТРАПЕЦИИ
% =========================================================

\section*{8. Трапеция: дополнительное построение через равнобедренный треугольник}
\topicline

Пусть $ABCD$ --- трапеция,

\[
AD\parallel BC.
\]

Точка $E$ --- середина боковой стороны $AB$,
причём

\[
EC=ED.
\]

Требуется доказать,
что трапеция прямоугольная.

Через $E$ проведём прямую

\[
EF\parallel AD,
\]

где

\[
F\in CD.
\]

Так как $AD\parallel BC$,
получаем также

\[
EF\parallel BC.
\]

Точка $E$ является серединой $AB$.
Прямая, проходящая через середину одной боковой стороны трапеции
параллельно основаниям,
проходит через середину второй боковой стороны.

Следовательно,

\[
CF=FD.
\]

По условию

\[
EC=ED.
\]

Значит, треугольник $ECD$ равнобедренный.

В нём $F$ --- середина основания $CD$,
поэтому $EF$ является медианой.

Медиана, проведённая к основанию равнобедренного треугольника,
одновременно является высотой:

\[
EF\perp CD.
\]

Но

\[
EF\parallel AD.
\]

Следовательно,

\[
AD\perp CD.
\]

Значит,

\[
\angle ADC=90^\circ.
\]

Следовательно,
трапеция $ABCD$ прямоугольная.

\begin{center}
\begin{tikzpicture}[scale=0.82]

    \tkzDefPoint(0,0){A}
    \tkzDefPoint(2,4){B}
    \tkzDefPoint(7,4){C}
    \tkzDefPoint(7,0){D}

    \tkzDefMidPoint(A,B)
    \tkzGetPoint{E}

    \tkzDefMidPoint(C,D)
    \tkzGetPoint{F}

    \tkzDrawPolygon(A,B,C,D)

    \tkzDrawSegments(
        E,C
        E,D
        E,F
    )

    \tkzDrawPoints(
        A,B,C,D,E,F
    )

    \tkzLabelPoints[below](A,D)
    \tkzLabelPoints[above](B,C)
    \tkzLabelPoints[left](E)
    \tkzLabelPoints[right](F)

    \tkzMarkSegments[mark=|](A,E E,B)
    \tkzMarkSegments[mark=||](C,F F,D)
    \tkzMarkSegments[mark=x](E,C E,D)

    \tkzMarkRightAngle[size=0.28](E,F,D)

\end{tikzpicture}
\end{center}

\textbf{Концепция поиска.}
Если в условии появляется равнобедренный треугольник,
сразу вспоминайте медиану, высоту и биссектрису,
проведённые к его основанию.

Дополнительное построение должно активировать
известное свойство фигуры.

% =========================================================
% 9. ТРАПЕЦИЯ И ПОДОБИЕ
% =========================================================

\section*{9. Трапеция: диагональ и подобные треугольники}
\topicline

Пусть в трапеции $ABCD$

\[
AD=45,
\qquad
BC=15.
\]

Точка $N$ лежит на боковой стороне $CD$,
причём

\[
CN=12,
\qquad
ND=18.
\]

Через точку $N$ проведена прямая

\[
KN\parallel AD\parallel BC,
\]

где

\[
K\in AB.
\]

Требуется найти $KN$.

Проведём диагональ $AC$ и обозначим

\[
Q=AC\cap KN.
\]

Треугольники

\[
\triangle CQN
\quad\text{и}\quad
\triangle CAD
\]

подобны.

Действительно,

\[
QN\parallel AD.
\]

Коэффициент подобия:

\[
\frac{CN}{CD}
=
\frac{12}{12+18}
=
\frac{12}{30}
=
\frac25.
\]

Поэтому

\[
\frac{QN}{AD}
=
\frac25.
\]

Отсюда

\[
QN=
45\cdot\frac25
=
18.
\]

Теперь рассмотрим треугольник $ABC$.

Так как

\[
KQ\parallel BC,
\]

то

\[
\triangle AKQ
\sim
\triangle ABC.
\]

Параллельные прямые
$AD$, $KN$ и $BC$
отсекают на боковых сторонах трапеции
пропорциональные отрезки.

Поэтому

\[
\frac{AK}{AB}
=
\frac{DN}{DC}
=
\frac{18}{30}
=
\frac35.
\]

Из подобия

\[
\frac{KQ}{BC}
=
\frac35.
\]

Следовательно,

\[
KQ=
15\cdot\frac35
=
9.
\]

Искомый отрезок:

\[
KN=
KQ+QN.
\]

Значит,

\[
\boxed{
KN=9+18=27
}.
\]

\begin{center}
\begin{tikzpicture}[scale=0.78]

    \tkzDefPoint(0,0){A}
    \tkzDefPoint(2,4){B}
    \tkzDefPoint(6,4){C}
    \tkzDefPoint(9,0){D}

    \tkzDefPoint(7.2,2.4){N}
    \tkzDefPoint(1.2,2.4){K}

    \tkzInterLL(A,C)(K,N)
    \tkzGetPoint{Q}

    \tkzDrawPolygon(A,B,C,D)

    \tkzDrawSegments(
        K,N
        A,C
    )

    \tkzDrawPoints(
        A,B,C,D,K,N,Q
    )

    \tkzLabelPoints[below](A,D)
    \tkzLabelPoints[above](B,C)
    \tkzLabelPoints[left](K)
    \tkzLabelPoints[right](N)
    \tkzLabelPoints[above left](Q)

    \tkzLabelSegment[below](A,D){$45$}
    \tkzLabelSegment[above](B,C){$15$}
    \tkzLabelSegment[right](C,N){$12$}
    \tkzLabelSegment[right](N,D){$18$}

\end{tikzpicture}
\end{center}

% =========================================================
% БЛОК-СХЕМА 3
% =========================================================

\clearpage

\section*{Алгоритм: поиск идеи в задаче на трапецию}

\vspace{8mm}

\begin{center}
\begin{tikzpicture}[node distance=8mm]

    \node[flowstart] (s)
    {Дана трапеция};

    \node[
        flowdecision,
        below=10mm of s
    ] (d1)
    {
        Есть ли в условии середина стороны
        или прямая, параллельная основаниям?
    };

    \node[
        flowaction,
        below left=13mm and 9mm of d1,
        text width=40mm
    ] (a1)
    {
        Проверьте среднюю линию
        и теоремы о пропорциональном
        делении боковых сторон
    };

    \node[
        flowaction,
        below right=13mm and 9mm of d1,
        text width=40mm
    ] (a2)
    {
        Проведите одну диагональ
        или высоту,
        чтобы получить треугольники
    };

    \node[
        flowaction,
        below=36mm of d1
    ] (a3)
    {
        Ищите пары подобных треугольников
        по параллельным прямым
    };

    \node[
        flowdecision,
        below=10mm of a3
    ] (d2)
    {
        Появился равнобедренный
        или прямоугольный треугольник?
    };

    \node[
        flowaction,
        below left=13mm and 9mm of d2,
        text width=40mm
    ] (a4)
    {
        Подключите его специальные свойства:
        медиану, высоту,
        теорему Пифагора, площадь
    };

    \node[
        flowaction,
        below right=13mm and 9mm of d2,
        text width=40mm
    ] (a5)
    {
        Используйте пропорции подобия
        и длины параллельных сечений
    };

    \node[
        flowstart,
        below=37mm of d2
    ] (f)
    {
        Получите искомую длину
        или требуемый угол
    };

    \draw[flowarrow]
    (s)--(d1);

    \draw[flowarrow]
    (d1)--node[above left,font=\small]{да}(a1);

    \draw[flowarrow]
    (d1)--node[above right,font=\small]{нет}(a2);

    \draw[flowarrow]
    (a1.south) |- (a3.west);

    \draw[flowarrow]
    (a2.south) |- (a3.east);

    \draw[flowarrow]
    (a3)--(d2);

    \draw[flowarrow]
    (d2)--node[above left,font=\small]{да}(a4);

    \draw[flowarrow]
    (d2)--node[above right,font=\small]{нет}(a5);

    \draw[flowarrow]
    (a4.south) |- (f.west);

    \draw[flowarrow]
    (a5.south) |- (f.east);

\end{tikzpicture}
\end{center}

\clearpage

% =========================================================
% 10. ТИПИЧНЫЕ ОШИБКИ
% =========================================================

\section*{10. Типичные ошибки и самопроверка}
\topicline

\begin{enumerate}

    \item
    \textbf{Потеря ОДЗ после сокращения.}

    Выколотые точки исходной рациональной функции
    остаются запрещёнными.

    \item
    \textbf{Включение нуля знаменателя
    в ответ неравенства.}

    Знаменатель никогда не может быть равен нулю.

    \item
    \textbf{Механическое построение параболы целиком.}

    После раскрытия модуля каждая формула действует
    только на своей части оси $x$.

    \item
    \textbf{Случайный выбор точек.}

    Для параболы полезны:
    вершина,
    пересечения с осями
    и симметричные пары точек.

    \item
    \textbf{Замена без ограничения.}

    При

    \[
    t=y^2
    \]

    обязательно фиксируется

    \[
    t\ge0.
    \]

    Если дополнительно известно $y\ne0$,
    то

    \[
    t>0.
    \]

    \item
    \textbf{Формула без происхождения.}

    Высоту к гипотенузе удобно помнить
    через равенство площадей:

    \[
    \frac12ab
    =
    \frac12ch.
    \]

    \item
    \textbf{Лишние построения в геометрии.}

    Сначала ищите треугольники.
    Диагональ, высота или прямая через середину стороны
    обычно дают нужную структуру.

    \item
    \textbf{Использование свойства,
    которое ещё не доказано.}

    Если требуется доказать прямоугольность трапеции,
    нельзя заранее использовать прямой угол как данность.

\end{enumerate}

% =========================================================
% ТРЕНИРОВОЧНЫЙ БЛОК
% =========================================================

\clearpage

\section*{Тренировочный блок --- Путь А}
\topicline

Задачи расположены по нарастающей сложности.
Решения в этом блоке не приводятся.

\begin{enumerate}[label=\textbf{\arabic*.}]

    \item
    Решите неравенство

    \[
    \frac{-18}{x^2+x-12}\le0.
    \]

    \item
    Два мастера выполняют одинаковый заказ
    объёмом $120$ деталей.

    Второй мастер изготавливает в час
    на $4$ детали меньше первого
    и поэтому тратит на заказ
    на $5$ часов больше.

    Составьте уравнение
    и найдите производительности обоих мастеров.

    \item
    В прямоугольном треугольнике
    гипотенуза равна $25$,
    один катет равен $15$.

    Найдите высоту,
    проведённую к гипотенузе.

    \item
    Решите систему

    \[
    \begin{cases}
    x^2+y^2=25,\\
    xy=-12.
    \end{cases}
    \]

    двумя способами:
    через подстановку
    и через формулу квадрата суммы или разности.

    \item
    Постройте график функции

    \[
    y=x^2+6|x|-7.
    \]

    Определите наибольшее число общих точек
    этого графика с прямой,
    параллельной оси $Ox$.

    \item
    В трапеции $ABCD$ основания

    \[
    AD=50,
    \qquad
    BC=20.
    \]

    Точка

    \[
    N\in CD,
    \]

    причём

    \[
    CN=9,
    \qquad
    ND=21.
    \]

    Через $N$ проведена прямая

    \[
    KN\parallel AD,
    \qquad
    K\in AB.
    \]

    Найдите $KN$.

    \item
    В трапеции $ABCD$ выполнено

    \[
    AD\parallel BC.
    \]

    Точка $E$ --- середина $AB$,
    а

    \[
    EC=ED.
    \]

    Через $E$ проведена прямая

    \[
    EF\parallel AD,
    \qquad
    F\in CD.
    \]

    Докажите, что

    \[
    EF\perp CD,
    \]

    а затем установите вид трапеции.

    \item
    Для функции

    \[
    y=
    \frac{(x-3)(x^2-1)}
         {x^2-2x-3}
    \]

    постройте график с учётом ОДЗ
    и найдите все значения $k$,
    при которых прямая

    \[
    y=kx
    \]

    не имеет с графиком общих точек.

    \item
    Решите неравенство

    \[
    \frac{(x-5)(x+3)}
         {(x+2)(x-1)}
    \ge0.
    \]

    В ответе отдельно отметьте,
    какие критические точки включаются
    и почему.

    \item
    Для функции

    \[
    y=
    \frac{(x-2)(x^2+4)}
         {x-2}
    \]

    исследуйте ОДЗ,
    постройте график
    и найдите все значения параметра $k$,
    при которых прямая

    \[
    y=kx
    \]

    не имеет с графиком общих точек.

\end{enumerate}



\end{document}